\documentclass[10pt,a4paper]{amsart}
\usepackage[margin=19mm]{geometry}
\usepackage{amsmath,amssymb,mathtools}
\usepackage[scr=rsfs,cal=euler]{mathalfa}
\usepackage{xfrac}
\usepackage[unicode,hidelinks,bookmarks=false]{hyperref}
\newcommand{\ie}{i.e.\ }
\newcommand{\cf}{cf.\ }
\newcommand{\resp}{respectively\ }
\newcommand{\Subsection}{\subsection}
\providecommand{\nobreakdash}{\nobreak\hbox{-}}
\setlength{\emergencystretch}{2em}
\multlinegap0pt
\usepackage{xcolor}
\newtheorem{proposition}{Proposition}
\usepackage{cleveref}
\crefname{proposition}{Proposition}{Propositions}
\newcommand*{\tran}{^{\mkern-1.5mu\mathsf{T}}} %Transpose
\title{Dictionary learning for Kernel EDMD}
\author{Erik Lien Bolager, Boumediene Hamzi, Houman Owhadi, Ioannis G. Kevrekidis, Felix Dietrich}
\date{Source: arXiv:2604.25572v1 (CC BY 4.0)}
\begin{document}
\maketitle

\noindent\emph{Source and licence.} Dictionary learning for Kernel EDMD. Version arXiv:2604.25572v1 (CC BY 4.0), distributed by arXiv under the Creative Commons Attribution 4.0 International licence, http://creativecommons.org/licenses/by/4.0/, as recorded in the arXiv licence metadata for this exact version and shown on the abstract page https://arxiv.org/abs/2604.25572v1. Full licence notice: \texttt{licenses/arxiv-2604.25572-CC-BY-4.0.txt}.\par\smallskip

\noindent\textit{Editorial scope.} Counted unit: the article's proposition prop:simp\_kedmd (source0.tex lines 251--261). The article's complete original proof of it is delivered with it (source0.tex lines 262--291, one \texttt{proof} environment of 30 lines). No mathematics is written, repaired or re-derived: the statement and the proof are the article's own lines, in the article's own notation, and no retained line is substituted or re-flowed.  Retained uncounted as the source-local context the proof needs: lines 132--134.  Nothing was omitted from any retained span, so the package carries no omission record.  No retained line contains a citation, so the wrapper carries no bibliography and no bibliography entry.  No retained line contains a figure environment, an image-inclusion command or drawing code, so no image is delivered and the package declares no asset dependency.  Editorial formatting only, in addition: an AMS \texttt{amsart} preamble that restates the article's own declarations and macros used by the retained lines, namely the environments \texttt{proposition} on separate counters, together with the standard packages those lines need.  The retained environments print in this document as Proposition 1 in order of appearance, whereas the article prints the same items with the numbering of its own full document.  The complete native source of the article as distributed in the arXiv e-print for this exact version is bundled unchanged as \texttt{sources/199311/source0.tex}.
\begin{align}\label{eq:def_khat}
    K_P = (P\Psi(Y))(P\Psi(X))^+ = Q(Z\Sigma^+)\tran ((P\Psi(X))\tran (P\Psi(Y))) (Z\Sigma^+) Q\tran = Q \Hat K Q\tran.
\end{align}

\begin{proposition}\label{prop:simp_kedmd}
    Let \(Z\) and \(\Sigma\) be defined as in \Cref{eq:def_khat}. There exists a linear map between \(K_{tr}\) and \(K_{sk}\), namely,
    \begin{align*}
        K_{tr} = (Z \Sigma)^+ K_{sk} Z\Sigma.
    \end{align*}
    We then have the following relations for the eigenvalue/eigenvector pairs of \(K_{tr}\) and \(K_{sk}\),
    \begin{itemize}
        \item If \((\lambda, w_{tr}, v_{tr})\) are the eigenvalue and corresponding left/right eigenvectors for \(K_{tr}\), then \((\lambda, w_{tr} \Sigma^+ Z\tran, Z\Sigma v_{tr})\) are the eigenvalue and corresponding left/right eigenvectors for \(K_{sk}\).
        \item If \((\lambda, w_{sk}, v_{sk})\) are the eigenvalue and corresponding left/right eigenvectors for \(K_{sk}\), then \((\lambda, w_{sk} Z\Sigma, \Sigma^+Z\tran v_{sk})\) are the eigenvalue and corresponding left/right eigenvectors for \(K_{tr}\).
    \end{itemize}
\end{proposition}

\begin{proof}
   First recall that \(g_{\theta}(X,X) = Z\Sigma^2 Z\tran\). We then have
   \begin{align*}
        K_{sk} &= g_{\theta}(Y,X) Z \Sigma^+ \Sigma^+ Z\tran = g_{\theta}(Y,X) (Z\Sigma^+) (Z \Sigma^+)\tran,
   \end{align*} 
   and by recalling from \Cref{eq:def_khat} that \(K_{tr} = (Z\Sigma^+)\tran g_{\theta}(Y,X) (Z\Sigma^+)\), we find
   \begin{align*}
        K_{sk} &=  (Z \Sigma) (\Sigma^+Z\tran) g_{\theta}(Y,X) (Z\Sigma^+) (Z \Sigma^+)\tran = (Z\Sigma) K_{tr} (Z\Sigma^+)\tran = (Z \Sigma)K_{tr} (Z\Sigma)^+.
   \end{align*}
   Equivalently, we have
   \begin{align*}
        K_{tr} = (Z\Sigma)^+ K_{sk} (Z\Sigma).
   \end{align*}
   For completeness we also show the relations for the eigenvalue and eigenvectors. Assume \((\lambda, w_{tr}, v_{tr})\) are eigenvalue and corresponding eigenvectors of \(K_{tr}\). Then 
   \begin{align*}
        w_{sk}K_{sk} = w_{tr} \Sigma^+ Z\tran K_{sk} = w_{tr} (Z\Sigma)^+ K_{sk} (Z\Sigma) (Z\Sigma)^+ = w_{tr} K_{tr} (Z\Sigma)^+ = \lambda w_{tr}(Z\Sigma)^+ = \lambda w_{sk},
   \end{align*}
   and
   \begin{align*}
        K_{sk} v_{sk} = K_{sk} Z\Sigma v_{tr} = (Z\Sigma) (Z\Sigma)^+ K_{sk} (Z\Sigma) v_{tr} = \lambda (Z\Sigma) v_{tr} = \lambda v_{sk}.
   \end{align*}
   Similarly, assuming \((\lambda, w_{sk}, v_{sk})\) are eigenvalue and corresponding eigenvectors of \(K_{sk}\), we have
   \begin{align*}
        w_{tr}K_{tr} = w_{sk} Z\Sigma K_{tr} = w_{sk} (Z\Sigma) K_{tr} (Z\Sigma)^+ (Z\Sigma) = \lambda w_{sk} Z\Sigma = \lambda w_{tr}
   \end{align*}
   and
   \begin{align*}
        K_{sk} v_{sk} = K_{sk} Z\Sigma v_{tr} = (Z\Sigma) (Z\Sigma)^+ K_{sk} (Z\Sigma) v_{tr} = \lambda (Z\Sigma) v_{tr} = \lambda v_{sk}.
   \end{align*}
\end{proof}

\end{document}
